\chapter{Categorical Irreversibility and the Arrow of Time}
\label{ch:arrow_of_time}

\begin{chapterabstract}
The arrow of time---the asymmetry between past and future in a time-symmetric physical law---is explained by categorical irreversibility: completed categorical states cannot be re-occupied. This is not the thermodynamic arrow (entropy increase), which is derived \emph{from} categorical irreversibility in Chapter~\ref{ch:triple_entropy}, but the categorical arrow that underlies it. The phenomenological experience of time flowing forward (Chapter~\ref{ch:time}) is the subjective face of categorical irreversibility: we experience the categorical arrow from the inside. The block universe is fully compatible with this picture: the mathematical structure is timeless, but its physical instantiation requires sequential tracing, and the direction of that tracing is fixed by categorical irreversibility.
\end{chapterabstract}

\section{The Standard Arrow of Time}

The standard explanation of the arrow of time appeals to the low-entropy initial conditions of the universe combined with the second law of thermodynamics. This explanation is circular: the second law itself requires explanation, and appealing to initial conditions merely pushes the question back.

\section{Categorical Irreversibility as Foundation}

Within the categorical framework, the arrow of time is not derived from entropy but the reverse: entropy (Chapter~\ref{ch:triple_entropy}) is derived from categorical irreversibility. The No Null State axiom and the completion sequence ordering together imply that categorical states can only be visited once, in order. This is the arrow. Entropy increase is its macroscopic statistical signature.

\section{The Phenomenological Arrow}

The subjective experience of time as flowing forward is the first-person report of categorical irreversibility. Each completed circuit cannot be re-entered; the sequence cannot be reversed; the past is inaccessible not because of thermodynamic cost but because categorical completion is a one-way operation. The ``specious present'' ($\sim$100--1,000~ms, Chapter~\ref{ch:time}) is the window of active circuit completion; the past is the accumulation of completed (and therefore inaccessible) categorical states.
